% !TEX program = lualatex
\documentclass[a4paper,11pt]{scrartcl}

\usepackage[
  lang=en,
  mode=journal,
  theme=graphite,
  fontset=lm,
  toc=true,
  toc-layout=page
]{synaptic}
\usepackage{booktabs,tabularx}
\newcolumntype{Y}{>{\raggedright\arraybackslash}X}

\SynapticHeader{Synaptic Technical Reference}
\SynapticSubHeader{Architecture, invariants, and release workflow}
\SynapticPreprint{v3.0.0}
\SynapticTitle{synaptic Technical Reference}
\SynapticAuthor{synaptic project}
\SynapticAffiliation{Open-source LaTeX community}
\SynapticSubject{Technical reference for synaptic version 3.0.0}
\SynapticKeywords{LaTeX3, architecture, l3build, regression testing}

\begin{abstract}
  This reference describes synaptic's module graph, load-time invariants,
  namespace rules, font and colour subsystems, extension points, generated-file
  policy, automated verification workflow, and packaging details.
\end{abstract}

\begin{document}
\SynapticMakeTitle

\section{Architecture}

\begin{verbatim}
synaptic.sty
  `-- synaptic-base.sty
      |-- synaptic-lang-{en,zh}.def
      |-- synaptic-color.sty
      |-- synaptic-fonts.sty
      |-- synaptic-layout.sty
      |-- synaptic-theorem.sty
      |-- synaptic-title.sty       shared metadata/renderers
      `-- mode module(s)
          |-- synaptic-journal.sty     journal
          |-- synaptic-book.sty        book
          |-- synaptic-lecture.sty     lecture
          `-- synaptic-notes.sty       notes
\end{verbatim}

The thin shell forwards user options to \texttt{synaptic-base} and defines a
small set of documentation helpers. The base module checks
the engine and class, parses options, loads labels, validates book mode, and
dispatches core and mode-specific modules. Language definition files are loaded
before any module so every label constant is available everywhere.

\subsection{Module responsibilities}

\noindent{\setlength{\parindent}{0pt}\noindent\begin{tabularx}{\textwidth}{@{}l l Y@{}}
  \toprule
  \textbf{Module} & \textbf{Kind} & \textbf{Responsibility} \\
  \midrule
  \texttt{synaptic.sty} & shell & option forwarding, documentation helpers \\
  \texttt{synaptic-base.sty} & core & engine/class checks, keys, dispatch, setup \\
  \texttt{synaptic-lang-*.def} & data & bilingual label constants \\
  \texttt{synaptic-color.sty} & core & theme registry, semantic tokens, delivery \\
  \texttt{synaptic-fonts.sty} & core & font bundles, auto selection, CJK fonts \\
  \texttt{synaptic-layout.sty} & core & geometry, typography, headings, numbering \\
  \texttt{synaptic-theorem.sty} & core & theorem styles, environments, proofs \\
  \texttt{synaptic-title.sty} & core & metadata model and four title renderers \\
  \texttt{synaptic-journal.sty} & mode & journal navigation and statements \\
  \texttt{synaptic-book.sty} & mode & chapter/part style, matter helpers \\
  \texttt{synaptic-lecture.sty} & mode & masthead, teaching cards, navigation \\
  \texttt{synaptic-notes.sty} & mode & notes masthead and knowledge cards \\
  \bottomrule
\end{tabularx}}

\noindent{\setlength{\parindent}{0pt}\noindent\begin{tabularx}{\textwidth}{@{}l l Y@{}}
  \toprule
  \textbf{Mode} & \textbf{Class invariant} & \textbf{Extra modules} \\
  \midrule
  journal & KOMA article/report & journal \\
  book & chapter-based KOMA class with front matter & book \\
  lecture & KOMA article/report & lecture \\
  notes & KOMA article/report & notes \\
  \bottomrule
\end{tabularx}}

\section{Configuration lifecycle}

The package-level \texttt{synaptic} key family is evaluated before module
loading. Mode, language, font set, title layout, measure, density, numbering,
delivery colour mode, extra math alphabets, and fine-breaking behaviour
therefore become structural invariants. The runtime key family
exposes only the state that can be applied without reloading packages:
\texttt{theme}, \texttt{toc}, and \texttt{toc-layout}. Structural or unknown
keys in a runtime request are errors; unknown package options are warnings.

Book mode is rejected unless both chapter and front-matter interfaces exist.
LuaLaTeX and KOMA-Script requirements are fatal diagnostics rather than delayed
undefined-control-sequence failures.

\section{Design invariants}

The following properties are deliberately asserted by the implementation and
protected by the regression suite.

\begin{itemize}
  \item \textbf{Deterministic typography.} Font auto-selection and theme
    derivation never depend on the machine beyond installed fonts; line
    spacing, measure, and heading sizes are explicit.
  \item \textbf{No global bad-box suppression.} Emergencystretch and tolerance
    are tuned; overfull and underfull warnings remain visible for the author.
  \item \textbf{No forced paper size or page-number reset.} Geometry follows
    the class; journal titles do not silently restart numbering.
  \item \textbf{One Unicode conversion.} PDF metadata is passed to hyperref as
    source tokens; pre-conversion would double-encode.
  \item \textbf{Collision-resistant public names.} Generic aliases are
    installed only when free.
  \item \textbf{Semantic colour architecture.} Components consume semantic
    tokens; themes and delivery modes change tokens, never component logic.
\end{itemize}

\section{Subsystems}

\subsection{Colour tokens}

Built-in themes are stored in a global property list. Activating a theme maps
its primary, secondary, and accent colours to semantic public names and derives
surface, border, rule, link, readable-muted, decorative-muted, and text tokens.
Print and mono delivery modes alter derived roles without changing component
code. Danger, warning, and success
colours remain semantic constants so their meaning does not change with brand
colour.

\begin{verbatim}
\SynapticDefineTheme{name}{primary}{secondary}[accent]
\SynapticUseTheme{name}
\end{verbatim}

The arguments are xcolor expressions. Defining a custom theme registers and
activates it.

\subsection{Font selection}

Font discovery is implemented entirely with \verb|\IfFontExistsTF|. Each
candidate is a complete roman/sans/mono/math bundle; auto selection never picks
a family whose companions are missing. The fixed priority is XCharter,
Libertinus, STIX Two, then Latin Modern. An explicit incomplete choice emits a
package error with the missing bundle requirements.

For \texttt{lang=zh}, ctex is loaded without a preset font scheme. Noto CJK SC
is preferred and Fandol is the TeX Live fallback. Synthetic weight/slant is
limited to faces for which the preferred family has no native counterpart.

\subsection{Layout and document structure}

Geometry combines mode defaults with explicit measure and density profiles.
Journal and lecture favour moderate line lengths; notes supports compact
single/two-column composition; book uses mirrored inner/outer margins, binding
correction, outer running heads, and empty duplex blanks. Widow and club penalties,
vertical rhythm, list spacing, line spacing, and paragraph indentation are
explicit rather than hidden behind global bad-box suppression.

KOMA heading fonts and skips are installed in the
\texttt{begindocument/before} hook, with all numbered levels sharing the sans
heading family: the primary colour for \texttt{section}, dark for
\texttt{subsection}, and a smaller, muted tint for \texttt{subsubsection} so
the depth hierarchy stays legible. Hyperref, bookmark, microtype, captions,
lists, caption punctuation, and mode-aware equation numbering are configured centrally.
A hairline running-head separator (\texttt{headsepline}) is installed by the
mode modules through a shared helper and coloured by the dedicated KOMA
\texttt{headsepline} font, keeping decoration lighter than head text.

The table of contents is refined through KOMA's \texttt{\string\DeclareTOCStyleEntry}
immediately before it is typeset: top-level entries repeat the on-page heading
colour in bold, deeper entries remain dark, and number alignment and vertical
rhythm are tightened. Figure and table caption labels carry the secondary
theme token.

\subsection{Titles and metadata}

Title values are held in expl3 token lists and sequences. Authors and their
affiliation marks use parallel sequences; affiliations and emails are
repeatable. A common metadata model feeds four renderers: inline journal,
book title sequence, course masthead, and compact notes masthead. The title
command validates a non-empty title and controls optional TOC placement without
silently resetting the user's page numbering. Standard LaTeX metadata is an
input adapter, not a redefinition of \verb|\maketitle|.

PDF information is passed as source tokens to hyperref. Hyperref performs the
Unicode PDF-string conversion once; pre-converting and then calling
\verb|\hypersetup| would double-encode metadata.

\subsection{Theorems}

thmtools defines plain, definition, and remark styles on top of amsthm. The
built-in statement family (\texttt{theorem}, \texttt{lemma},
\texttt{proposition}, \texttt{definition}, \texttt{corollary}, \texttt{axiom},
\texttt{conjecture}, \texttt{claim}, \texttt{property}, \texttt{example},
\texttt{exercise}, \texttt{criterion}, \texttt{convention},
\texttt{problem}, \texttt{solution}, \texttt{remark}, \texttt{note}, and
\texttt{warning}) uses a mode-aware shared counter: chapter scope for books and
section scope elsewhere by default. Under \texttt{theorem-style=boxed} the
numbered statement environments are wrapped with tcolorbox (a theme-tinted
background with a west accent bar) at \texttt{\string\begin{document}}; proofs
stay unboxed. \texttt{theorem-style=plain} disables that wrapping. Examples and
exercises are plain-style statements; warnings use the remark presentation.
Because every statement shares the theorem counter, the module rebuilds each
environment's hyperlink anchor (\texttt{theH<name>}) against the real structural
parent so links stay unique across sections and chapters.

\section{Namespace policy}

\noindent{\setlength{\parindent}{0pt}\noindent\begin{tabularx}{\textwidth}{@{}l l Y@{}}
  \toprule
  \textbf{Scope} & \textbf{Pattern} & \textbf{Example} \\
  \midrule
  Public command & \texttt{\textbackslash Synaptic\ldots} & \SynapticCmd{SynapticDefineTheme} \\
  Public environment & \texttt{syn...} & \texttt{synsummary} \\
  Internal function & \texttt{\textbackslash\_\_synaptic\_\ldots:} & title rendering helper \\
  Module interface & \texttt{\textbackslash synaptic\_\ldots:} & mode accessor \\
  Local state & \texttt{\textbackslash l\_synaptic\_\ldots} & selected theme token list \\
  Global registry & \texttt{\textbackslash g\_synaptic\_\ldots} & theme property list \\
  Label constant & \texttt{\textbackslash c\_synaptic\_label\_\ldots} & theorem label \\
  \bottomrule
\end{tabularx}}

New public names should avoid common nouns. Statement environments are part of
the built-in family: declare new ones with \verb|\SynapticNewTheorem| rather
than adding aliases.

\section{Extension points}

\begin{itemize}
  \item Add labels in a new \texttt{synaptic-lang-*.def} file, then extend the
    language key and definition-file dispatch.
  \item Add a complete font predicate, setter, explicit dispatch branch, and a
    deliberate location in the auto-selection chain.
  \item Add modes in the base key family, layout dispatch, and mode-module
    loader. State the required KOMA class capability and test it.
  \item Build visual extensions from semantic colour tokens rather than fixed
    RGB values.
  \item Add tests along with every extension; the suite is the guard for the
    invariants above.
\end{itemize}

\section{Automated verification}

\texttt{synaptic.dtx} is the source of truth. DocStrip uses
\texttt{synaptic.ins} to generate the installable \texttt{.sty} and
\texttt{.def} files. The generated copies under
\texttt{tex/latex/synaptic/} are committed for zero-install use and must remain
byte-identical to fresh extraction.

The regression suite is organised as follows:

{\setlength{\parindent}{0pt}\noindent\begin{tabularx}{\textwidth}{@{}l Y@{}}
  \toprule
  \textbf{Group} & \textbf{Coverage} \\
  \midrule
  \texttt{journal}, \texttt{lecture}, \texttt{notes}, \texttt{book} & core
  mode behaviour, navigation, and defaults \\
  \texttt{book-zh}, \texttt{lecture-zh}, \texttt{journal-zh},
  \texttt{language-zh} & Chinese labels, CJK fonts, and bilingual layering \\
  \texttt{themes} & cycling all five themes and a custom theme \\
  \texttt{color-mode-print}, \texttt{options} & print and monochrome delivery \\
  \texttt{customization} & runtime setup, custom theorems \\
  \texttt{toc-layouts}, \texttt{title-layouts} & TOC and title compositions \\
  \texttt{layout-profiles} & density, measure, binding offset \\
  \texttt{book-full} & complete title sequence and matter workflow \\
  \texttt{lecture-continuous} & continuous numbering profile \\
  \texttt{extra-math} & optional Unicode math alphabets \\
  \bottomrule
\end{tabularx}}

Each group typesets a document with the standard engine and compares the
normalised log against its committed expectation file. CI runs the suite and
additionally builds the documented source, four manuals, and all examples,
failing on warnings, overfull boxes, or missing characters.

\section{Source and release workflow}

\begin{verbatim}
l3build unpack     # regenerate .sty/.def into build/local
l3build check      # run the regression suite
l3build doc        # build the documented source, manuals
l3build ctan       # create TDS and CTAN archives
l3build install   # copy package files into the texmf tree
l3build uninstall # remove a previous installation
\end{verbatim}

The \texttt{ctan} target emits \texttt{synaptic-ctan.zip} (the source tree
plus generated documentation and examples) together with the nested
\texttt{synaptic.tds.zip} in TDS layout. Installation replaces the previous
version in place; because the generated file set is stable,
\texttt{l3build install} is safe to rerun after every release.

\section{Migrating to 3.0}

Version 3.0 removes every backward-compatibility layer: the \verb|\pkg| and
\verb|\cmd| aliases, the lowercase book matter helpers, the import of standard
\verb|\title|/\verb|\author|/\verb|\date| metadata, the \texttt{box-numbering}
option, and the \texttt{synaptic-boxes} module. Examples, warnings, and
exercises are now first-class members of the built-in statement family;
\texttt{title-layout=cover} became \texttt{title-layout=sequence}; re-declaring
an environment with \verb|\SynapticNewTheorem| is a hard error. New documents
should use the namespaced public forms exclusively.

\end{document}
